\section{Conclusions}\label{sec:conclusions}

Our goal was to measure how much a land-use classifier trained on EuroSAT loses in semi-arid
Aguascalientes, and how much of that loss light fine-tuning on local labels wins back. We
compared ResNet-50, EfficientNet-B0 and ViT-S/16 under one shared protocol, and three findings
stand out.

\begin{enumerate}
  \item \emph{The models were almost tied on EuroSAT, but all of them dropped sharply in
  Aguascalientes (RQ-a, RQ-b).} EuroSAT macro-F1 ranged from \EuroSATMin{} to \EuroSATMax, and
  without local training the drop ranged from \DropMin{} to \DropMax. This means that a
  near-perfect EuroSAT score does not guarantee a useful model in another region; the model has
  to be evaluated on local labels.
  \item \emph{Light fine-tuning won back only part of the drop (RQ-c).} Training the classifier
  alone (\texttt{head}) recovered at most \HeadRecoveryMax{} of the drop, and our runs cannot
  tell whether the features were too weak or the classifier was undertrained. Fine-tuning the
  last block (\texttt{last\_block}) recovered between \RecoveryMin{} and \RecoveryMax{}
  (Eq.~\ref{eq:recovery}). This means that most of the drop remained in every run after
  fine-tuning.
  \item \emph{The effect of the input bands and the weights depended on the model (RQ-d,
  RQ-e).} With ImageNet weights, 13 bands helped only ResNet-50 in Aguascalientes. With 13
  bands, the SSL4EO-S12 runs gave the highest observed macro-F1 after \texttt{last\_block} for
  both ResNet-50 and ViT-S, but they also used another normalization, so the gap cannot be
  credited to the weights alone. This means that the answer depends on the combination of model, bands and weights, not
  only on CNN versus ViT.
\end{enumerate}

The selection rule, fixed before any result, chose ResNet-50/MS13/SSL4EO with
\texttt{last\_block} fine-tuning, with a macro-F1 of \AgsLastBlockBest, \SelectionMargin{} above
the runner-up. Among the six ImageNet runs, where the three models compete on equal weights,
ViT-S/RGB/ImageNet came first with \ImageNetLastBlockBest, \ImageNetSelectionMargin{} above the
next run. On cost (RQ-a), EfficientNet-B0 had far fewer parameters but a similar time per epoch,
and ViT-S/RGB/ImageNet had the highest observed throughput, \PatchesPerSecMax{} patches/s.

What we learned is that the domain shift outweighed the differences between the models, so a
model for a new region has to be judged on local labels and a spatially separated test set, not
on EuroSAT alone. The selected model is the starting point for the second stage of the
project, a monthly land-use monitor of the state. Section~\ref{sec:discussion} lists the
experiments that would make these answers firmer.
